\chapter{算法不变量、有限随机性与证明分层}

\section{先把随机性变成显式输入}

普通程序常在函数内部调用随机数生成器，这不利于数学推理。形式化时可以把所有随机选择看成显式输入 $r$：
\[
\operatorname{run}:\mathrm{Input}\to\mathrm{RandomTape}\to\mathrm{Output}.
\]
固定 $r$ 后，\code{run x r} 是确定函数；概率来自对 random tape 的分布。这样可以把问题拆成两层：
\begin{enumerate}
  \item 对每条随机带证明状态安全、输出类型正确等确定性性质；
  \item 对随机带集合计数或使用概率测度，证明成功事件占多大比例。
\end{enumerate}

\begin{intuitionbox}
随机算法不是“不可预测的代码”，而是“确定程序加一个按指定分布采样的额外输入”。这个视角让执行语义与概率分析分离，也便于用有限枚举、Lean 和外部测试交叉验证。
\end{intuitionbox}

\section{状态机、不变量与后置条件}

流算法可写成：
\[
s_0=\operatorname{init},\qquad
s_{t+1}=\operatorname{step}(s_t,x_{t+1},r_{t+1}),\qquad
y=\operatorname{query}(s_n).
\]
至少区分三类性质：
\begin{description}
  \item[表示不变量] 状态字段长度、索引、非负性和内部一致性；
  \item[语义不变量] 状态和已处理输入前缀之间的关系；
  \item[最终保证] 在终止或查询时由不变量推出输出正确性/误差界。
\end{description}

上一章的 \code{StreamState.Inv} 只保证列表长度不超过已见元素数，它是表示不变量，不说明保留的元素具有正确抽样分布。

\begin{warningbox}
“状态大小不超过 $k$”不是 reservoir sampling 正确性的全部。真正分布性结论是处理 $n$ 个元素后，每个元素以 $k/n$ 概率进入样本（在相应条件下），还需要形式化随机选择及其分布。
\end{warningbox}

\section{有限均匀样本空间}

入门时可以绕开测度论，令有限类型 $\Omega$ 表示所有随机带，均匀概率定义为
\[
\Pr[E]=\frac{|\{\omega\in\Omega:E(\omega)\}|}{|\Omega|},
\qquad |\Omega|>0.
\]
Lean 中可先定义事件计数：
\begin{lstlisting}
def eventCount [Fintype Omega]
    (E : Omega -> Prop) [DecidablePred E] : Nat :=
  (Finset.univ.filter E).card

def uniformProb [Fintype Omega]
    (E : Omega -> Prop) [DecidablePred E] : Rat :=
  (eventCount E : Rat) / (Fintype.card Omega : Rat)
\end{lstlisting}
若 $\Omega$ 为空，分母为零，必须在定理假设中排除或改用带非空证明的样本空间结构。定义“能通过类型检查”不等于数学接口已完整。

\section{indicator 与期望}

事件 $E$ 的指示变量
\[
\mathbf{1}_E(\omega)=
\begin{cases}
1,&E(\omega),\\
0,&\text{otherwise}
\end{cases}
\]
满足 $\mathbb{E}[\mathbf{1}_E]=\Pr[E]$。有限均匀空间中，这一事实归结为：对所有 $\omega$ 求和时，每个满足 $E$ 的元素贡献 1，其余贡献 0，因此总和恰为事件基数。

形式化路线可以拆成：
\begin{enumerate}
  \item 定义可判定事件和 indicator；
  \item 证明 indicator 的有限和等于 filter 后的 card；
  \item 除以非零的 $|\Omega|$；
  \item 再把结果包装成期望/概率接口。
\end{enumerate}
这比直接挑战一般概率库中的积分定理更适合作为第一题。

\section{有限 union bound 的计数骨架}

对事件 $E_1,\ldots,E_m$，union bound 为
\[
\Pr\!\left[\bigcup_{i=1}^m E_i\right]
\le \sum_{i=1}^m \Pr[E_i].
\]
有限均匀空间下只需证明集合基数不等式
\[
\left|\bigcup_i E_i\right|\le\sum_i |E_i|,
\]
再除以 $|\Omega|$。重叠元素在右边可能被重复计数，所以得到上界。Lean 实现时要选择 \code{Finset}、\code{Set} 或有限索引族，并处理 membership、decidability 与 coercion。

\section{随机算法保证的六层拆分}

\begin{center}
\begin{tabularx}{0.96\textwidth}{>{\bfseries}lX X}
\toprule
层 & 典型命题 & 合适证据 \\
\midrule
接口 & 输入可解析、输出类型正确 & 类型检查、编译 \\
状态安全 & 索引合法、不变量保持 & Lean 归纳定理 \\
逐随机带正确性 & 对所有 $r$ 都满足确定性关系 & Lean theorem \\
分布正确性 & random tape 确实服从所需分布 & 有限计数或概率库定理 \\
概率保证 & 失败概率不超过 $\delta$ & 期望、尾界、union bound 的形式证明 \\
资源保证 & 时间、空间、随机位成本 & 成本语义、复杂度证明或审计分析 \\
\bottomrule
\end{tabularx}
\end{center}

一个项目可以只形式化其中一部分，但必须标明其余层是测试、纸笔证明、外部假设还是完全未处理。

\section{Count-Min 的分层形式化示例}

Count-Min Sketch 的端到端结论可以先拆成：
\begin{enumerate}
  \item 数组更新保持所有计数器非负；
  \item 目标项每行对应桶包含 $f_i$，所以逐行估计不小于 $f_i$；
  \item 碰撞噪声是其他频率的加权和；
  \item 通用哈希给出碰撞概率，推出噪声期望界；
  \item Markov 与独立重复给出高概率误差界；
  \item 最小值查询和参数 $w,d$ 合成最终 theorem。
\end{enumerate}
第一、二步基本不含概率，适合先做；后三步需要哈希族和有限概率接口。完整形式化不是“一条巨大 theorem”，而是一条依赖图。

\section{实现与规范的连接}

Lean 中容易证明的纯函数可能与 Python/C++ 高性能实现不同。至少有三种连接策略：
\begin{enumerate}
  \item 以 Lean 定义为可执行 reference implementation，外部实现做 differential testing；
  \item 为外部语言建立语义模型，再证明实现 refinement；
  \item 只证明抽象算法 theorem，同时明确工程实现尚未形式化。
\end{enumerate}
第一种最适合个人学习项目。它不能自动证明外部实现正确，但能给测试提供可信 oracle。

\begin{aibox}
AI 生成的随机算法候选应先通过分层 verifier：接口检查、小规模枚举、hidden random tests、确定性不变量证明、概率保证审查和复杂度审查。Lean 最适合其中可精确定义的层，不应被当作覆盖所有研究判断的万能标签。
\end{aibox}

\section{建议的第一个随机算法形式化题}

选择如下梯度：
\begin{enumerate}
  \item 有限集合上 indicator 的和等于事件基数；
  \item 均匀概率总和为 1（显式假设样本空间非空）；
  \item 两个或有限多个事件的 union bound；
  \item 一个抽样状态的大小不变量；
  \item 简化抽样器的无偏性；
  \item 最后才进入完整 reservoir sampling、Chernoff 或 Count-Min。
\end{enumerate}

\begin{checkpointbox}
\begin{enumerate}
  \item 把随机算法写成输入和 random tape 上的确定函数；
  \item 区分表示不变量、语义不变量和分布正确性；
  \item 在有限均匀空间中推导 indicator 期望等于事件概率；
  \item 写出 union bound 的集合计数证明骨架；
  \item 把 Count-Min 端到端保证拆成至少五条 lemma；
  \item 为一个只完成部分形式化的项目写清证据边界。
\end{enumerate}
\end{checkpointbox}
